\documentclass[11pt]{article}
\usepackage[margin=1in]{geometry}
\usepackage{amsmath,amssymb}
\usepackage[hidelinks]{hyperref}

\title{Cubic Integration of Well-Rounded Logconcave Functions:\\
A Later-Literature Answer to arXiv:1409.6011}
\author{Run fa\_banach\_001, agent\_lane\_16}
\date{2026-08-11}

\begin{document}
\maketitle

\section*{Status}

\textbf{Literature already answered (affirmative), with limited scope.}
Conclusion item 2 of Cousins--Vempala is answered by Theorem 1.8 of
Kook--Vempala.  The separate cubic arbitrary-body rounding question in
Conclusion item 4 remains open.

\section*{Source question}

Ben Cousins and Santosh Vempala, arXiv:1409.6011, ask in Conclusion item 2 on
PDF page 35:
\begin{quote}
It would be interesting to extend our algorithm to integrating any
well-rounded logconcave function; we expect this should be possible with
essentially the same complexity.
\end{quote}
The preceding item records the relevant dimension complexity for their
well-rounded convex-body volume algorithm as $O^*(n^3)$.

\section*{Explicit answering theorem}

Yunbum Kook and Santosh S. Vempala, arXiv:2411.13462, explicitly identify
this provenance: in their technical overview they write that
Cousins--Vempala raised the question of extending the annealing strategy to
arbitrary logconcave distributions, and their Result 4 says that they extend
the Cousins--Vempala annealing approach to general logconcave integration.

Their Theorem 1.8 on PDF page 7 states that, for every $\varepsilon>0$ and
every integrable well-rounded logconcave function
$f=e^{-V}:\mathbb R^n\to\mathbb R$ presented by the specified evaluation
oracle, a randomized algorithm returns a $(1+\varepsilon)$-multiplicative
approximation to $\int f$ with probability at least $3/4$, using
\[
  \widetilde O\!\left(\frac{n^3}{\varepsilon^2}\right)
\]
evaluation queries.  The theorem also gives
$\widetilde O(n^{3.5}\operatorname{polylog}R+n^3/\varepsilon^2)$ for an
arbitrary, not initially well-rounded, logconcave function.

This is exactly the requested extension in the well-rounded case: the
dimension exponent is cubic, with the usual accuracy dependence made
explicit.  Because the answering authors cite and describe the original
question, the relationship is an explicit literature answer rather than an
agent-inferred theorem match.

\section*{Scope limitation}

This packet closes only Conclusion item 2.  Conclusion item 4 asks for an
$O^*(n^3)$ rounding algorithm for arbitrary convex bodies.  The current 2026
revision of Jia--Laddha--Lee--Vempala, arXiv:2008.02146, gives
$\widetilde O(n^{3.5})$ with the current KLS bound, and the 2026 survey section
of Kook--Vempala, arXiv:2507.18021, still calls this the best known general
rounding complexity.  No cubic general rounding algorithm is claimed here.

\section*{Search evidence}

The identification was checked against the source text of arXiv:1409.6011,
the current source of arXiv:2008.02146, the source and official PDF of
arXiv:2411.13462, and the 2026 discussion in arXiv:2507.18021.  The exact
answering statement is Theorem 1.8 of arXiv:2411.13462; its overview explicitly
cites the source work and describes the extension.

\begin{thebibliography}{9}
\bibitem{CV}
B. Cousins and S. S. Vempala,
\emph{Gaussian Cooling and $O^*(n^3)$ Algorithms for Volume and Gaussian
Volume}, arXiv:1409.6011.

\bibitem{KV}
Y. Kook and S. S. Vempala,
\emph{Sampling and Integration of Logconcave Functions by Algorithmic
Diffusion}, arXiv:2411.13462, Theorem 1.8.

\bibitem{JLLV}
H. Jia, A. Laddha, Y. T. Lee, and S. S. Vempala,
\emph{Reducing Isotropy and Volume to KLS: Faster Rounding and Volume
Algorithms}, arXiv:2008.02146.
\end{thebibliography}

\end{document}
